% contributions.tex -- research questions and contributions (draft 4 Oct 2026). Follows the
% introduction. The research questions are those of the project plan; the contributions list
% what has been built and shown so far, and is to be trimmed to what the final results support.

\subsection{Research Questions}
\label{sec:rq}

The research question of this thesis is whether a centralised controller can manage the join and
leave transients of individual agents in a multi-UAV payload system while keeping the payload
stable. Two questions follow from it.
\begin{enumerate}[label=\textbf{RQ\arabic*}]
    \item Can a centralised planner that has the whole team's state redistribute the load when a UAV
    leaves the payload and when a new UAV attaches to it mid-flight, with a transient small enough
    for the transport to continue?
    \item Does a decentralised dissipative control law \cite{quan2025swarm} or a trajectory-based
    optimal control approach \cite{tudelft} better balance robustness and tracking performance
    for this problem?
\end{enumerate}
The working hypothesis is that the centralised approach suits a change of team composition,
because the coordinating node can recompute the force distribution for the new team at the moment
of the event, rather than relying on the remaining agents to absorb it passively.

\subsection{Contributions}
\label{sec:contributions}

The work to date contributes the following.
\begin{enumerate}
    \item \textbf{A centralised payload controller flown on commercial-grade hardware.} The
    load--cable planner of \cite{tudelft} is combined with a cable-aware tracking MPC per UAV that
    commands Betaflight flight controllers over their radio links, with no changes to the
    UAVs' firmware. It includes a floor start for rigid rods, a pretension and lift procedure, and
    a thrust map identified on the rig (Sections~\ref{sec:controller} and~\ref{sec:takeoff}).
    \item \textbf{Reconfiguration by resizing the optimal control problem.} The attachment
    geometry is a runtime parameter of the planner, so a detach or an attach becomes a switch to
    a pre-built solver for the new team size, with each UAV kept on its real attachment point
    (Section~\ref{sec:reconfig}).
    \redo{OCP resize against handing the payload to a dissipative network during the transient, on a detach and on an attach.}
% --- original text, flown before the rig twin; kept to refill this place:
%     (Section~\ref{sec:reconfig}). In simulation this outperforms handing the payload to a
%     dissipative network during the transient.
    \item \textbf{A geometric condition for a level detach.} Whether the payload can stay level
    after a UAV leaves depends on whether its centre of mass stays inside the polygon of the
    surviving attachment points. This sets the attachment layout for four UAVs.
    \item \textbf{Mid-flight attachment of a new UAV.} A free UAV welds onto the carried payload
    through an electromagnet and is folded into the planner, in combination with a collaborator's
    approach controller.
    \redo{the full sequence in simulation: a UAV leaves, completes a pick-up task and rejoins while the team flies an orbit (the rejoin mission).}
% --- original text, flown before the rig twin; kept to refill this place:
%     approach controller. The full sequence (a UAV leaves, completes a pick-up task and rejoins
%     while the team flies an orbit) is demonstrated in simulation.
    \item \textbf{A comparison of the trajectory-based and dissipative approaches} on the same
    plant, tracker and test cases, including the tracking architecture the dissipative method
    assumes.
    \redo{the comparison itself (Section~\ref{sec:sim_architecture}).}
    \item \textbf{An analysis of steady errors in the cascade.} Several sources of steady payload
    error were identified and separated: a relaxed initial condition in the tracker, a breakaway
    pull cap held for the whole flight, model errors in thrust and rod mass, and external forces.
    Integral action in the planner's model and several estimation approaches from the
    literature are compared against them.
    \item \textbf{A simulation and verification environment that mirrors the rig.} A Gazebo twin
    of the lab set-up runs the same launch files and parameters as the hardware. Every run is
    recorded against a named baseline, with a set number of repeats.
    \item \textbf{Hardware results} with three UAVs carrying a 0.86\,kg ring: hover at up to
    1.0\,m, a circle and a figure-eight, and an account of the gap between simulation and the rig.
\end{enumerate}
